\documentclass[conference,compsoc]{IEEEtran}
\usepackage{longtable}
\usepackage{booktabs}
\usepackage{array}
\usepackage{textcomp}
\usepackage{needspace}

\begin{document}
\onecolumn

\section{Table Definitions}

% Suppress warnings about missing aux file on first compile
\IfFileExists{\jobname.aux}{\tableofcontents\clearpage}{}

\needspace{4\baselineskip}
\subsection{all\_files}
\label{sec:all_files}

Master list of filenames seen in the collect. Derived from process executions, dll loads and file usage activity.

\begin{longtable}{|p{5cm}|p{2cm}|p{9cm}|}
\hline
\textbf{Column} & \textbf{Data Type} & \textbf{Description} \\
\hline
\endfirsthead

\multicolumn{2}{c}%
{\tablename\ \thetable\ -- Continued from previous page} \\
\hline
\textbf{Column} & \textbf{Data Type} & \textbf{Description} \\
\hline
\endhead

\hline
\multicolumn{2}{r}{Continued on next page} \\
\endfoot

\hline
\endlastfoot

dll\_num\_rows & DOUBLE & Number of times used as a DLL \\
file\_num\_rows & DOUBLE & Number of times referenced in file read/write activity \\
filename & STRING & Fully qualified filename \\
num\_hosts & INT64 & Number of unique values in the related column. Typically, the related column should only have 1 value. \\
process\_num\_rows & DOUBLE & Number of rows file was executed as a process \\
\hline
\end{longtable}

\needspace{4\baselineskip}
\subsection{files}
\label{sec:files}

Files by hostname with summary information for executions, dlls and file usage.

\begin{longtable}{|p{5cm}|p{2cm}|p{9cm}|}
\hline
\textbf{Column} & \textbf{Data Type} & \textbf{Description} \\
\hline
\endfirsthead

\multicolumn{2}{c}%
{\tablename\ \thetable\ -- Continued from previous page} \\
\hline
\textbf{Column} & \textbf{Data Type} & \textbf{Description} \\
\hline
\endhead

\hline
\multicolumn{2}{r}{Continued on next page} \\
\endfoot

\hline
\endlastfoot

dll\_first\_seen & TIMESTAMP & Earliest time seen \\
dll\_last\_seen & TIMESTAMP & Latest time seen \\
dll\_num\_rows & DOUBLE & Number of times used as a DLL \\
file\_first\_seen & TIMESTAMP & Earliest time seen \\
file\_id & STRING & Unique ID for a file. Hash of hostname+fully qualified filename \\
file\_last\_seen & TIMESTAMP & Latest time seen \\
file\_num\_rows & DOUBLE & Number of times referenced in file read/write activity \\
filename & STRING & Fully qualified filename \\
hostname & STRING & Hostname data collected from \\
max\_process\_term & TIMESTAMP & Maximum value in the interval \\
min\_process\_started & TIMESTAMP & Minimum value in the interval \\
process\_num\_rows & DOUBLE & Number of rows file was executed as a process \\
\hline
\end{longtable}

\needspace{4\baselineskip}
\subsection{host}
\label{sec:host}

Unique entry for every instrumented host.

\begin{longtable}{|p{5cm}|p{2cm}|p{9cm}|}
\hline
\textbf{Column} & \textbf{Data Type} & \textbf{Description} \\
\hline
\endfirsthead

\multicolumn{2}{c}%
{\tablename\ \thetable\ -- Continued from previous page} \\
\hline
\textbf{Column} & \textbf{Data Type} & \textbf{Description} \\
\hline
\endhead

\hline
\multicolumn{2}{r}{Continued on next page} \\
\endfoot

\hline
\endlastfoot

Hostname & STRING & Hostname data collected from \\
ad\_domain & STRING & Active directory host is joined to \\
agent\_ids & STRING[] & List of agent ids. A long running host may have several dur to update, reinstalls, etc. \\
arch & STRING & Chip architecture \\
domain\_role & STRING & Host role within the domain. Ex: domain controller, member, etc. \\
etl\_version & STRING & Wintap data collection plugin version \\
first\_seen & TIMESTAMP & Earliest time seen \\
has\_battery & BOOLEAN & Does this host have a battery? Indicates laptop if true \\
last\_boot & INT64 & Time of last boot, in epoch seconds \\
last\_seen & TIMESTAMP & Latest time seen \\
num\_ad\_domain & INT64 & Number of unique values in the related column. Typically, the related column should only have 1 value. \\
num\_arch & INT64 & Number of unique values in the related column. Typically, the related column should only have 1 value. \\
num\_domain\_role & INT64 & Number of unique values in the related column. Typically, the related column should only have 1 value. \\
num\_etl\_version & INT64 & Number of unique values in the related column. Typically, the related column should only have 1 value. \\
num\_has\_battery & INT64 & Number of unique values in the related column. Typically, the related column should only have 1 value. \\
num\_last\_boot & INT64 & Number of unique values in the related column. Typically, the related column should only have 1 value. \\
num\_os & INT64 & Number of unique values in the related column. Typically, the related column should only have 1 value. \\
num\_os\_version & INT64 & Number of unique values in the related column. Typically, the related column should only have 1 value. \\
num\_processor\_count & INT64 & Number of unique values in the related column. Typically, the related column should only have 1 value. \\
num\_processor\_speed & INT64 & Number of unique values in the related column. Typically, the related column should only have 1 value. \\
num\_rows & INT64 & Number of row behind this summary \\
num\_wintap\_version & INT64 & Number of unique values in the related column. Typically, the related column should only have 1 value. \\
os & STRING & Host operating system \\
os\_family & STRING & OS Family: windows, linux, osx \\
os\_version & STRING & Operating System Version \\
processor\_count & INTEGER & Number of processors cores \\
processor\_speed & INT64 & Speed of processors \\
wintap\_version & STRING & Version Wintap core. Each plugin has its own version. \\
\hline
\end{longtable}

\needspace{4\baselineskip}
\subsection{host\_ip}
\label{sec:host_ip}

Network interface information for hosts.

\begin{longtable}{|p{5cm}|p{2cm}|p{9cm}|}
\hline
\textbf{Column} & \textbf{Data Type} & \textbf{Description} \\
\hline
\endfirsthead

\multicolumn{2}{c}%
{\tablename\ \thetable\ -- Continued from previous page} \\
\hline
\textbf{Column} & \textbf{Data Type} & \textbf{Description} \\
\hline
\endhead

\hline
\multicolumn{2}{r}{Continued on next page} \\
\endfoot

\hline
\endlastfoot

Hostname & STRING & Hostname data collected from \\
MTU & INTEGER & Should be max transmition unit, but its currently not collected. \\
agent\_id & STRING & Agent ID is a unique id for a Wintap install on a host. It is created on initial startup after install. \\
first\_seen & TIMESTAMP & Earliest time seen \\
interface & STRING & Should be the interface device name, but its currently not collected. \\
ip\_addr & INT64 & An IP address on the host. Stored as an integer \\
ip\_addr\_no & STRING & An IP address on the host. In dotted quad notation. \\
last\_seen & TIMESTAMP & Latest time seen \\
mac & STRING & Mac address as string (colon delimited) \\
num\_rows & INT64 & Number of row behind this summary \\
os\_family & STRING & OS Family: windows, linux, osx \\
private\_gateway & STRING & IP of the gateway used by this interface \\
\hline
\end{longtable}

\needspace{4\baselineskip}
\subsection{labels\_graph\_net\_conn}
\label{sec:labels_graph_net_conn}

Summarizes network-related labels by conn\_id

\begin{longtable}{|p{5cm}|p{2cm}|p{9cm}|}
\hline
\textbf{Column} & \textbf{Data Type} & \textbf{Description} \\
\hline
\endfirsthead

\multicolumn{2}{c}%
{\tablename\ \thetable\ -- Continued from previous page} \\
\hline
\textbf{Column} & \textbf{Data Type} & \textbf{Description} \\
\hline
\endhead

\hline
\multicolumn{2}{r}{Continued on next page} \\
\endfoot

\hline
\endlastfoot

conn\_id & STRING & Hash of "normalized" 5 tuple (ip1, port1, ip2, port2, L4 protocol). To represent A to B the same as B to A, the lower IP (as int) along with its port comes first in the hash preimage ordering. When ip1 == ip2, then the pair order is decided by lowest port instead. \\
label\_num\_hits & INT64 & Number of unique labels hit \\
label\_num\_sources & INT64 & Number of unique sources of labels \\
label\_num\_uniq\_annotations & INT64 & Number of unqiue annotations \\
label\_source & STRING & For now, only a single source: "networkx" \\
\hline
\end{longtable}

\needspace{4\baselineskip}
\subsection{labels\_graph\_nodes}
\label{sec:labels_graph_nodes}

Parse JSON in NetworkX files into a flat structure representing the nodes

\begin{longtable}{|p{5cm}|p{2cm}|p{9cm}|}
\hline
\textbf{Column} & \textbf{Data Type} & \textbf{Description} \\
\hline
\endfirsthead

\multicolumn{2}{c}%
{\tablename\ \thetable\ -- Continued from previous page} \\
\hline
\textbf{Column} & \textbf{Data Type} & \textbf{Description} \\
\hline
\endhead

\hline
\multicolumn{2}{r}{Continued on next page} \\
\endfoot

\hline
\endlastfoot

annotation & STRING & Annotations are manually added to specific nodes by the person creating the networkx graph of activity. Used to further highlight the activity. \\
filename & STRING & Fully qualified filename \\
id & STRING & Unique ID for the node in the graph. Often a PID\_HASH, CONN\_ID or similar. \\
label & STRING & Graph node label. Often a PROCESS\_NAME or similar. \\
node\_type & STRING & Type of node: PROCESS, NETWORK, FILE, etc. \\
\hline
\end{longtable}

\needspace{4\baselineskip}
\subsection{labels\_graph\_process\_summary}
\label{sec:labels_graph_process_summary}

Summarizes process-related labels by pid\_hash

\begin{longtable}{|p{5cm}|p{2cm}|p{9cm}|}
\hline
\textbf{Column} & \textbf{Data Type} & \textbf{Description} \\
\hline
\endfirsthead

\multicolumn{2}{c}%
{\tablename\ \thetable\ -- Continued from previous page} \\
\hline
\textbf{Column} & \textbf{Data Type} & \textbf{Description} \\
\hline
\endhead

\hline
\multicolumn{2}{r}{Continued on next page} \\
\endfoot

\hline
\endlastfoot

label\_num\_hits & INT64 & Number of unique labels hit \\
label\_num\_sources & INT64 & Number of unique sources of labels \\
label\_num\_uniq\_annotations & INT64 & Number of unqiue annotations \\
label\_source & STRING & For now, only a single source: "networkx" \\
pid\_hash & STRING & Globally unique process hash (hash from hostname, os pid, and start\_time) \\
\hline
\end{longtable}

\needspace{4\baselineskip}
\subsection{labels\_networkx}
\label{sec:labels_networkx}

Import of networkx graph data. Used primarily as base to build summary tables. Data in this table is JSON, tricky to work with directly.

\begin{longtable}{|p{5cm}|p{2cm}|p{9cm}|}
\hline
\textbf{Column} & \textbf{Data Type} & \textbf{Description} \\
\hline
\endfirsthead

\multicolumn{2}{c}%
{\tablename\ \thetable\ -- Continued from previous page} \\
\hline
\textbf{Column} & \textbf{Data Type} & \textbf{Description} \\
\hline
\endhead

\hline
\multicolumn{2}{r}{Continued on next page} \\
\endfoot

\hline
\endlastfoot

directed & BOOLEAN & Directed graph or not. \\
filename & STRING & Fully qualified filename \\
is\_multigraph & BOOLEAN & Multigraph or not \\
links & STRUCT & List of links \\
nodes & STRUCT & List of nodes \\
\hline
\end{longtable}

\needspace{4\baselineskip}
\subsection{lolbas}
\label{sec:lolbas}

Joins process data with LOLBAS (Living Off the Land Binaries and Scripts) data. Summarizes LOLBAS-related activity by pid\_hash.

\begin{longtable}{|p{5cm}|p{2cm}|p{9cm}|}
\hline
\textbf{Column} & \textbf{Data Type} & \textbf{Description} \\
\hline
\endfirsthead

\multicolumn{2}{c}%
{\tablename\ \thetable\ -- Continued from previous page} \\
\hline
\textbf{Column} & \textbf{Data Type} & \textbf{Description} \\
\hline
\endhead

\hline
\multicolumn{2}{r}{Continued on next page} \\
\endfoot

\hline
\endlastfoot

acknowledgements & STRING & Person who submitted the entry. May be more than 1. \\
author & STRING & Person or group who submitted the entry. \\
command & STRING & The command with arguments \\
command\_category & STRING & Intent. Ex: EXECUTE, UAC BYPASS, CONCEAL, etc. \\
command\_description & STRING & Description of the command \\
command\_privileges & STRING & Required priviledges to execute \\
command\_usecase & STRING & Description of use case for attacker \\
date & DATE & Date created \\
description & STRING & Normal purpose of the command \\
detections & STRING & How to detect exploit: Many are URLs pointing to SIGMA rule, Defender, etc. \\
filename & STRING & Filename of executable. No path. \\
mitre\_attck\_technique & STRING & Specific MITRE attck code. Ex: T1003 \\
operating\_system & STRING & OSes binary is found on. Some very specific, some general. \\
paths & STRING & Paths where binary is found. More than one are comma separated. \\
resources & STRING & URL to more info \\
url & STRING & URL to LolBAS project about the exploit \\
\hline
\end{longtable}

\needspace{4\baselineskip}
\subsection{mitre\_labels}
\label{sec:mitre_labels}

\begin{longtable}{|p{5cm}|p{2cm}|p{9cm}|}
\hline
\textbf{Column} & \textbf{Data Type} & \textbf{Description} \\
\hline
\endfirsthead

\multicolumn{2}{c}%
{\tablename\ \thetable\ -- Continued from previous page} \\
\hline
\textbf{Column} & \textbf{Data Type} & \textbf{Description} \\
\hline
\endhead

\hline
\multicolumn{2}{r}{Continued on next page} \\
\endfoot

\hline
\endlastfoot

analytic\_id & STRING & Mitre analytic ID. Ex: CAR-2021-05-012 \\
entity & STRING & Wintap entity ID (hash). The type of entity (PID\_HASH, CONN\_ID, etc) is defined in the ENTITY\_TYPE column. Note: Currently, only PID\_HASHes are identified. \\
entity\_type & STRING & Wintap entity type. Determines what kind of hash is in the ENTITY column. Ex: PID\_HASH, CONN\_ID, FILE\_ID \\
time & TIMESTAMP & Timestamp of activity \\
\hline
\end{longtable}

\needspace{4\baselineskip}
\subsection{process}
\label{sec:process}

Each row is a process execution. Processes are the central object of this entire data model.

\begin{longtable}{|p{5cm}|p{2cm}|p{9cm}|}
\hline
\textbf{Column} & \textbf{Data Type} & \textbf{Description} \\
\hline
\endfirsthead

\multicolumn{2}{c}%
{\tablename\ \thetable\ -- Continued from previous page} \\
\hline
\textbf{Column} & \textbf{Data Type} & \textbf{Description} \\
\hline
\endhead

\hline
\multicolumn{2}{r}{Continued on next page} \\
\endfoot

\hline
\endlastfoot

agent\_id & STRING & Agent ID is a unique id for a Wintap install on a host. It is created on initial startup after install. \\
args & STRING & Command line arguments \\
commit\_charge & INT64 & Commit charge is the total amount of virtual memory guaranteed for all processes to fit in physical memory and the page file. \\
commit\_peak & INT64 & Total memory commited \\
cpu\_cycle\_count & INT64 & Always zero (note: more than zero cycles were generally used in a process) \\
cpu\_utilization & INTEGER & Ignore: unable to collect. This *should* be a value indicating the percentage of cpu utilized by the process over its life, but turns out be harder to collect than expected. \\
exit\_code & INT64 & Process exit code \\
file\_id & STRING & Unique ID for a file. Hash of hostname+fully qualified filename \\
file\_md5 & STRING & MD5 hash of the binary process/dll \\
file\_sha2 & STRING & SHA2 hash of the file contents \\
filename & STRING & Fully qualified filename \\
first\_seen & TIMESTAMP & Earliest time seen \\
hard\_fault\_count & INTEGER & Count of hard page faults \\
hostname & STRING & Hostname data collected from \\
last\_seen & TIMESTAMP & Latest time seen \\
num\_agent\_id & INT64 & Number of unique values in the related column. Typically, the related column should only have 1 value. \\
num\_args & INT64 & Number of unique values in the related column. Typically, the related column should only have 1 value. \\
num\_file\_md5 & INT64 & Number of unique values in the related column. Typically, the related column should only have 1 value. \\
num\_file\_sha2 & INT64 & Number of unique values in the related column. Typically, the related column should only have 1 value. \\
num\_parent\_os\_pid & INT64 & Number of unique values in the related column. Typically, the related column should only have 1 value. \\
num\_parent\_pid\_hash & INT64 & Number of unique values in the related column. Typically, the related column should only have 1 value. \\
num\_process\_name & INT64 & Number of unique values in the related column. Typically, the related column should only have 1 value. \\
num\_process\_path & INT64 & Number of unique values in the related column. Typically, the related column should only have 1 value. \\
num\_process\_start & DOUBLE & Number of unique values in the related column. Typically, the related column should only have 1 value. \\
num\_process\_stop & DOUBLE & Number of unique values in the related column. Typically, the related column should only have 1 value. \\
num\_user\_name & INT64 & Number of unique values in the related column. Typically, the related column should only have 1 value. \\
os\_family & STRING & OS Family: windows, linux, osx \\
os\_pid & INTEGER & Operating system process id \\
parent\_os\_pid & INTEGER & OS Process ID of parent process \\
parent\_pid\_hash & STRING & PID\_HASH of parent process \\
pid\_hash & STRING & Globally unique process hash (hash from hostname, os pid, and start\_time) \\
process\_name & STRING & Filename of executable \\
process\_path & STRING & Fully qualified path of the process executable \\
process\_started & TIMESTAMP & Process start timestamp \\
process\_started\_seconds & DOUBLE & Process start time in unix epoch seconds \\
process\_stop\_seconds & DOUBLE & Process stop in unix epoch seconds \\
process\_term & TIMESTAMP & Process termination timestamp \\
read\_operation\_count & INT64 & Count of number of read operations \\
read\_transfer\_kilobytes & INT64 & Number of K bytes read \\
token\_elevation\_type & INTEGER & Indicates UAC privilege status: 1 (Default) = no linked token, UAC disabled or non-admin account; 2 (Full) = elevated administrator with full rights after UAC approval; 3 (Limited) = filtered admin token before elevation with available full token; 4 = undocumented value, treat as unknown/reserved. \\
user\_name & STRING & Username. Note that on windows there are many usernames that represent system or other background activity \\
write\_operation\_count & INT64 & Count of number of write operations \\
write\_transfer\_kilobytes & INT64 & Number of K bytes written \\
\hline
\end{longtable}

\needspace{4\baselineskip}
\subsection{process\_conn\_incr}
\label{sec:process_conn_incr}

Network connection increments for each process. (KEY: PID\_HASH + CONN\_ID + INCR\_START\_SECS)

\begin{longtable}{|p{5cm}|p{2cm}|p{9cm}|}
\hline
\textbf{Column} & \textbf{Data Type} & \textbf{Description} \\
\hline
\endfirsthead

\multicolumn{2}{c}%
{\tablename\ \thetable\ -- Continued from previous page} \\
\hline
\textbf{Column} & \textbf{Data Type} & \textbf{Description} \\
\hline
\endhead

\hline
\multicolumn{2}{r}{Continued on next page} \\
\endfoot

\hline
\endlastfoot

Hostname & STRING & Hostname data collected from \\
agent\_id & STRING & Agent ID is a unique id for a Wintap install on a host. It is created on initial startup after install. \\
conn\_id & STRING & Hash of "normalized" 5 tuple (ip1, port1, ip2, port2, L4 protocol). To represent A to B the same as B to A, the lower IP (as int) along with its port comes first in the hash preimage ordering. When ip1 == ip2, then the pair order is decided by lowest port instead. \\
first\_seen & TIMESTAMP & Earliest time seen \\
incr\_start & TIMESTAMP & Start time of a 1 minute increment used to aggregate the high-volume events \\
last\_seen & TIMESTAMP & Latest time seen \\
local\_ip\_addr & STRING & IP address on the host collecting data. This address is local to the host sensor. In dotted quad notation. \\
local\_ip\_int & INT64 & IP of the host collecting this data as 32-bit int \\
local\_port & INTEGER & Port of the process on the collecting host \\
max\_10sec\_eventcount & INTEGER & Maximum value in the interval \\
max\_size & INT64 & Maximum value in the interval \\
max\_tcp\_recv\_count & INTEGER & Maximum value in the interval \\
max\_tcp\_recv\_size & INT64 & Maximum value in the interval \\
max\_tcp\_send\_count & INTEGER & Maximum value in the interval \\
max\_tcp\_send\_size & INT64 & Maximum value in the interval \\
max\_udp\_recv\_count & INTEGER & Maximum value in the interval \\
max\_udp\_recv\_size & INT64 & Maximum value in the interval \\
max\_udp\_send\_count & INTEGER & Maximum value in the interval \\
max\_udp\_send\_size & INT64 & Maximum value in the interval \\
min\_10sec\_eventcount & INTEGER & Minimum value in the interval \\
min\_size & INT64 & Minimum value in the interval \\
min\_tcp\_recv\_size & INT64 & Minimum value in the interval \\
min\_tcp\_send\_size & INT64 & Minimum value in the interval \\
min\_udp\_recv\_size & INT64 & Minimum value in the interval \\
min\_udp\_send\_size & INT64 & Minimum value in the interval \\
num\_raw\_rows & INT64 & Number of unique values in the related column. Typically, the related column should only have 1 value. \\
os\_family & STRING & OS Family: windows, linux, osx \\
pid\_hash & STRING & Globally unique process hash (hash from hostname, os pid, and start\_time) \\
process\_name & STRING & Filename of executable \\
protocol & STRING & Layer 4 protocol (TCP, UDP) \\
remote\_ip\_addr & STRING & IP address the host collecting data is talking to. This address is remote to the host sensor. In dotted quad notation. \\
remote\_ip\_int & INT64 & IP of the remote host for this connection increment as 32-bit int \\
remote\_port & INTEGER & Port of the remote host for this connection increment \\
sq\_size & DOUBLE & Squared size of bytes in network packet. Ask Chris B! \\
sq\_tcp\_recv\_size & DOUBLE & Square of the values in the interval. (Ask Chris B) \\
sq\_tcp\_send\_size & DOUBLE & Square of the values in the interval. (Ask Chris B) \\
sq\_udp\_recv\_size & DOUBLE & Square of the values in the interval. (Ask Chris B) \\
sq\_udp\_send\_size & DOUBLE & Square of the values in the interval. (Ask Chris B) \\
tcp\_accept\_count & DOUBLE & Number of TCP ACCEPT events on this connection for the time window \\
tcp\_connect\_count & DOUBLE & Number of TCP CONNECT events on this connection for the time window \\
tcp\_disconnect\_count & DOUBLE & Number of TCP DISCONNECT events on this connection for the time window \\
tcp\_reconnect\_count & DOUBLE & Number of TCP RECONNECT events on this connection for the time window \\
tcp\_recv\_count & DOUBLE & Number of TCP RECV events on this connection for the time window \\
tcp\_recv\_size & DOUBLE & Number of TCP RECV bytes received on this connection for the time window \\
tcp\_retransmit\_count & DOUBLE & Number of TCP RETRANSMIT events on this connection for the time window \\
tcp\_send\_count & DOUBLE & Number of TCP SEND events on this connection for the time window \\
tcp\_send\_size & DOUBLE & Number of TCP bytes sent for the time window \\
tcp\_tcpcopy\_count & DOUBLE & Number of TCP TCPCOPY events on this connection for the time window \\
tcp\_tcpcopy\_size & DOUBLE & Number of TCP TCPCOPY bytes on this connection for the time window \\
total\_events & DOUBLE & Sum of events counts for this connection increment \\
total\_size & DOUBLE & Number of bytes observed for this connection increment \\
udp\_recv\_count & DOUBLE & Number of UDP RECV events on this connection for the time window \\
udp\_recv\_size & DOUBLE & Number of UDP bytes received for the time window \\
udp\_send\_count & DOUBLE & Number of UDP SEND events on this connection for the time window \\
udp\_send\_size & DOUBLE & Number of UDP bytes sent for the time window \\
\hline
\end{longtable}

\needspace{4\baselineskip}
\subsection{process\_exe\_file\_summary}
\label{sec:process_exe_file_summary}

Summary of files used in process executions. Derived from all process executions resulting in a single row per file per host.

\begin{longtable}{|p{5cm}|p{2cm}|p{9cm}|}
\hline
\textbf{Column} & \textbf{Data Type} & \textbf{Description} \\
\hline
\endfirsthead

\multicolumn{2}{c}%
{\tablename\ \thetable\ -- Continued from previous page} \\
\hline
\textbf{Column} & \textbf{Data Type} & \textbf{Description} \\
\hline
\endhead

\hline
\multicolumn{2}{r}{Continued on next page} \\
\endfoot

\hline
\endlastfoot

file\_id & STRING & Unique ID for a file. Hash of hostname+fully qualified filename \\
filename & STRING & Fully qualified filename \\
hostname & STRING & Hostname data collected from \\
max\_process\_term & TIMESTAMP & Maximum value in the interval \\
min\_process\_started & TIMESTAMP & Minimum value in the interval \\
process\_num\_rows & INT64 & Number of rows file was executed as a process \\
source & STRING & Source of file: currently only PROCESS \\
\hline
\end{longtable}

\needspace{4\baselineskip}
\subsection{process\_file}
\label{sec:process_file}

File activity of processes. Activity is summarized to the PROCESS + FILE level. See RAW\_PROCESS\_FILE for activity detailed activity.

\begin{longtable}{|p{5cm}|p{2cm}|p{9cm}|}
\hline
\textbf{Column} & \textbf{Data Type} & \textbf{Description} \\
\hline
\endfirsthead

\multicolumn{2}{c}%
{\tablename\ \thetable\ -- Continued from previous page} \\
\hline
\textbf{Column} & \textbf{Data Type} & \textbf{Description} \\
\hline
\endhead

\hline
\multicolumn{2}{r}{Continued on next page} \\
\endfoot

\hline
\endlastfoot

Hostname & STRING & Hostname data collected from \\
activity\_type & STRING & File activity type. Ex: READ, WRITE \\
agent\_id & STRING & Agent ID is a unique id for a Wintap install on a host. It is created on initial startup after install. \\
bytes\_requested & DOUBLE & Total bytes requested \\
event\_count & DOUBLE & Total ETW events \\
file\_hash & STRING & MD5 hash of the file \\
file\_id & STRING & Unique ID for a file. Hash of hostname+fully qualified filename \\
filename & STRING & Name of the file \\
first\_seen & TIMESTAMP & Earliest time seen \\
last\_seen & TIMESTAMP & Latest time seen \\
max\_event & TIMESTAMP & Maximum value in the interval \\
min\_event & TIMESTAMP & Minimum value in the interval \\
num\_raw\_rows & INT64 & Number of unique values in the related column. Typically, the related column should only have 1 value. \\
pid\_hash & STRING & Globally unique process hash (hash from hostname, os pid, and start\_time) \\
process\_name & STRING & Filename of executable \\
\hline
\end{longtable}

\needspace{4\baselineskip}
\subsection{process\_file\_summary}
\label{sec:process_file_summary}

File activity summarize to the process level.

\begin{longtable}{|p{5cm}|p{2cm}|p{9cm}|}
\hline
\textbf{Column} & \textbf{Data Type} & \textbf{Description} \\
\hline
\endfirsthead

\multicolumn{2}{c}%
{\tablename\ \thetable\ -- Continued from previous page} \\
\hline
\textbf{Column} & \textbf{Data Type} & \textbf{Description} \\
\hline
\endhead

\hline
\multicolumn{2}{r}{Continued on next page} \\
\endfoot

\hline
\endlastfoot

Close\_Events & DOUBLE & Number of file close events \\
Create\_Events & DOUBLE & Number of file create events \\
Delete\_Events & DOUBLE & Number of file delete events \\
Hostname & STRING & Hostname data collected from \\
Read\_Bytes & DOUBLE & File bytes read \\
Read\_Events & DOUBLE & Number of file read events \\
Rename\_Events & DOUBLE & Number of file rename events \\
SetInfo\_Events & DOUBLE &  \\
Write\_Bytes & DOUBLE & File bytes written \\
Write\_Events & DOUBLE & Number of file write events \\
agent\_id & STRING & Agent ID is a unique id for a Wintap install on a host. It is created on initial startup after install. \\
first\_seen & TIMESTAMP & Earliest time seen \\
last\_seen & TIMESTAMP & Latest time seen \\
num\_null\_filename & DOUBLE & Number of unique values in the related column. Typically, the related column should only have 1 value. \\
num\_raw\_rows & DOUBLE & Number of unique values in the related column. Typically, the related column should only have 1 value. \\
num\_uniq\_file\_hash & INT64 & Number of unique values in the related column. Typically, the related column should only have 1 value. \\
pid\_hash & STRING & Globally unique process hash (hash from hostname, os pid, and start\_time) \\
process\_name & STRING & Filename of executable \\
\hline
\end{longtable}

\needspace{4\baselineskip}
\subsection{process\_image\_load}
\label{sec:process_image_load}

DLLs loaded and unloaded by process.

\begin{longtable}{|p{5cm}|p{2cm}|p{9cm}|}
\hline
\textbf{Column} & \textbf{Data Type} & \textbf{Description} \\
\hline
\endfirsthead

\multicolumn{2}{c}%
{\tablename\ \thetable\ -- Continued from previous page} \\
\hline
\textbf{Column} & \textbf{Data Type} & \textbf{Description} \\
\hline
\endhead

\hline
\multicolumn{2}{r}{Continued on next page} \\
\endfoot

\hline
\endlastfoot

agent\_id & STRING & Agent ID is a unique id for a Wintap install on a host. It is created on initial startup after install. \\
build\_time & INT64 & Build time of DLL \\
checksum & INTEGER & A checksum on the file contents of the DLL \\
default\_base & STRING & Default base memory address \\
file\_id & STRING & Unique ID for a file. Hash of hostname+fully qualified filename \\
file\_md5 & STRING & MD5 hash of the binary process/dll \\
filename & STRING & Filename of the loaded code \\
first\_seen & TIMESTAMP & Earliest time seen \\
hostname & STRING & Hostname data collected from \\
image\_base & STRING & Actual base memory address(?) \\
last\_seen & TIMESTAMP & Latest time seen \\
max\_image\_size & INTEGER & Maximum value in the interval \\
min\_image\_size & INTEGER & Minimum value in the interval \\
num\_load & DOUBLE & Number of unique values in the related column. Typically, the related column should only have 1 value. \\
num\_unload & DOUBLE & Number of unique values in the related column. Typically, the related column should only have 1 value. \\
pid\_hash & STRING & Globally unique process hash (hash from hostname, os pid, and start\_time) \\
process\_name & STRING & Filename of executable \\
\hline
\end{longtable}

\needspace{4\baselineskip}
\subsection{process\_image\_load\_summary}
\label{sec:process_image_load_summary}

Summarizes DLL and image load activity by process (pid\_hash).

\begin{longtable}{|p{5cm}|p{2cm}|p{9cm}|}
\hline
\textbf{Column} & \textbf{Data Type} & \textbf{Description} \\
\hline
\endfirsthead

\multicolumn{2}{c}%
{\tablename\ \thetable\ -- Continued from previous page} \\
\hline
\textbf{Column} & \textbf{Data Type} & \textbf{Description} \\
\hline
\endhead

\hline
\multicolumn{2}{r}{Continued on next page} \\
\endfoot

\hline
\endlastfoot

agent\_id & STRING & Agent ID is a unique id for a Wintap install on a host. It is created on initial startup after install. \\
dlls & STRING[] & List of DDL names as an array \\
first\_seen & TIMESTAMP & Earliest time seen \\
hostname & STRING & Hostname data collected from \\
last\_seen & TIMESTAMP & Latest time seen \\
num\_uniq\_files & INT64 & Number of unique values in the related column. Typically, the related column should only have 1 value. \\
pid\_hash & STRING & Globally unique process hash (hash from hostname, os pid, and start\_time) \\
process\_name & STRING & Filename of executable \\
\hline
\end{longtable}

\needspace{4\baselineskip}
\subsection{process\_lolbas\_summary}
\label{sec:process_lolbas_summary}

\begin{longtable}{|p{5cm}|p{2cm}|p{9cm}|}
\hline
\textbf{Column} & \textbf{Data Type} & \textbf{Description} \\
\hline
\endfirsthead

\multicolumn{2}{c}%
{\tablename\ \thetable\ -- Continued from previous page} \\
\hline
\textbf{Column} & \textbf{Data Type} & \textbf{Description} \\
\hline
\endhead

\hline
\multicolumn{2}{r}{Continued on next page} \\
\endfoot

\hline
\endlastfoot

lolbas\_cats & STRING[] & List categories from hits \\
lolbas\_mitre & STRING[] & List of Mitre codes from hits \\
lolbas\_num\_rows & INT64 & Number of hits \\
lolbas\_privs & STRING[] & List of priviledges from hist \\
pid\_hash & STRING & Globally unique process hash (hash from hostname, os pid, and start\_time) \\
\hline
\end{longtable}

\needspace{4\baselineskip}
\subsection{process\_mitre\_summary}
\label{sec:process_mitre_summary}

\begin{longtable}{|p{5cm}|p{2cm}|p{9cm}|}
\hline
\textbf{Column} & \textbf{Data Type} & \textbf{Description} \\
\hline
\endfirsthead

\multicolumn{2}{c}%
{\tablename\ \thetable\ -- Continued from previous page} \\
\hline
\textbf{Column} & \textbf{Data Type} & \textbf{Description} \\
\hline
\endhead

\hline
\multicolumn{2}{r}{Continued on next page} \\
\endfoot

\hline
\endlastfoot

mitre\_analytic\_ids & STRING[] & List of analytic IDs from hits \\
mitre\_analytic\_types & STRING[][] & List of ananlytic types from hits \\
mitre\_information\_domains & STRING[] & List domains from hits. Ex: Analytic, Host, Network \\
mitre\_num\_rows & INT64 & Number of hits \\
mitre\_subtypes & STRING[][] & List of subtypes from hits. Ex: Map building, Anomaly, Hostflow.  Process. \\
pid\_hash & STRING & Globally unique process hash (hash from hostname, os pid, and start\_time) \\
\hline
\end{longtable}

\needspace{4\baselineskip}
\subsection{process\_net\_conn}
\label{sec:process_net_conn}

Network connections for each process. These are summarize to the PROCESS and Network 5-tuple (KEY: PID\_HASH + CONN\_ID)

\begin{longtable}{|p{5cm}|p{2cm}|p{9cm}|}
\hline
\textbf{Column} & \textbf{Data Type} & \textbf{Description} \\
\hline
\endfirsthead

\multicolumn{2}{c}%
{\tablename\ \thetable\ -- Continued from previous page} \\
\hline
\textbf{Column} & \textbf{Data Type} & \textbf{Description} \\
\hline
\endhead

\hline
\multicolumn{2}{r}{Continued on next page} \\
\endfoot

\hline
\endlastfoot

Hostname & STRING & Hostname data collected from \\
agent\_id & STRING & Agent ID is a unique id for a Wintap install on a host. It is created on initial startup after install. \\
conn\_id & STRING & Hash of "normalized" 5 tuple (ip1, port1, ip2, port2, L4 protocol). To represent A to B the same as B to A, the lower IP (as int) along with its port comes first in the hash preimage ordering. When ip1 == ip2, then the pair order is decided by lowest port instead. \\
first\_seen & TIMESTAMP & Earliest time seen \\
last\_seen & TIMESTAMP & Latest time seen \\
local\_ip\_addr & STRING & IP address on the host collecting data. This address is local to the host sensor. In dotted quad notation. \\
local\_port & INTEGER & Port of the process on the collecting host \\
num\_raw\_rows & DOUBLE & Number of unique values in the related column. Typically, the related column should only have 1 value. \\
os\_family & STRING & OS Family: windows, linux, osx \\
pid\_hash & STRING & Globally unique process hash (hash from hostname, os pid, and start\_time) \\
process\_name & STRING & Filename of executable \\
protocol & STRING & Layer 4 Protocol (TCP, UDP) \\
remote\_ip\_addr & STRING & IP address the host collecting data is talking to. This address is remote to the host sensor. In dotted quad notation. \\
remote\_port & INTEGER & Port of the remote host for this connection increment \\
sq\_size & DOUBLE & Squared size of bytes in network packet. Ask Chris B! \\
sq\_tcp\_recv\_size & DOUBLE & Square of the values in the interval. (Ask Chris B) \\
sq\_tcp\_send\_size & DOUBLE & Square of the values in the interval. (Ask Chris B) \\
sq\_udp\_recv\_size & DOUBLE & Square of the values in the interval. (Ask Chris B) \\
sq\_udp\_send\_size & DOUBLE & Square of the values in the interval. (Ask Chris B) \\
tcp\_accept\_count & DOUBLE & Number of TCP ACCEPT events on this connection \\
tcp\_connect\_count & DOUBLE & Number of TCP CONNECT events on this connection \\
tcp\_disconnect\_count & DOUBLE & Number of TCP DISCONNECT events on this connection \\
tcp\_reconnect\_count & DOUBLE & Number of TCP RECONNECT events on this connection \\
tcp\_recv\_count & DOUBLE & Number of TCP RECV events on this connection \\
tcp\_recv\_size & DOUBLE & Number of TCP RECV bytes received on this connection \\
tcp\_retransmit\_count & DOUBLE & Number of TCP RETRANSMIT events on this connection \\
tcp\_send\_count & DOUBLE & Number of TCP SEND events on this connection \\
tcp\_send\_size & DOUBLE & Number of TCP bytes sent \\
tcp\_tcpcopy\_count & DOUBLE & Number of TCP TCPCOPY events on this connection \\
tcp\_tcpcopy\_size & DOUBLE & Number of TCP TCPCOPY bytes on this connection \\
total\_events & DOUBLE & Sum of events counts for this connection \\
total\_size & DOUBLE & Number of bytes observed for this connection \\
udp\_recv\_count & DOUBLE & Number of UDP RECV events on this connection \\
udp\_recv\_size & DOUBLE & Number of UDP bytes received \\
udp\_send\_count & DOUBLE & Number of UDP SEND events on this connection \\
udp\_send\_size & DOUBLE & Number of UDP bytes sent on this connection \\
\hline
\end{longtable}

\needspace{4\baselineskip}
\subsection{process\_net\_summary}
\label{sec:process_net_summary}

Summarizes all network connections for a given process to the process level.

\begin{longtable}{|p{5cm}|p{2cm}|p{9cm}|}
\hline
\textbf{Column} & \textbf{Data Type} & \textbf{Description} \\
\hline
\endfirsthead

\multicolumn{2}{c}%
{\tablename\ \thetable\ -- Continued from previous page} \\
\hline
\textbf{Column} & \textbf{Data Type} & \textbf{Description} \\
\hline
\endhead

\hline
\multicolumn{2}{r}{Continued on next page} \\
\endfoot

\hline
\endlastfoot

Hostname & STRING & Hostname data collected from \\
agent\_id & STRING & Agent ID is a unique id for a Wintap install on a host. It is created on initial startup after install. \\
avg\_bytes & DOUBLE & Average bytes in an event \\
avg\_packets & DOUBLE & Average packets per session done by the process \\
conn\_id\_count & INT64 & Number of unique network connection IDs (5-tuples) \\
first\_seen & TIMESTAMP & Earliest time seen \\
last\_seen & TIMESTAMP & Latest time seen \\
max\_bytes & DOUBLE & Maximum bytes in an event \\
max\_packets & DOUBLE & Max packets per session done by the process \\
min\_bytes & DOUBLE & Minimum bytes in an event \\
min\_packets & DOUBLE & Min packets per session done by the process \\
net\_recv\_size & DOUBLE & Total bytes received by the process \\
net\_rs\_total & DOUBLE & Total bytes sent/received by the process \\
net\_send\_size & DOUBLE & Total bytes sent by the process \\
net\_send\_vs\_recv & DOUBLE & Ratio of bytes sent vs received by the process \\
net\_total\_events & DOUBLE & Total events accross UDP/TCP \\
net\_total\_size & DOUBLE & Total bytes accross all events \\
num\_raw\_rows & DOUBLE & Number of unique values in the related column. Typically, the related column should only have 1 value. \\
os\_family & STRING & OS Family: windows, linux, osx \\
pid\_hash & STRING & Globally unique process hash (hash from hostname, os pid, and start\_time) \\
process\_name & STRING & Filename of executable \\
sq\_size & DOUBLE & Squared size of bytes in network packet. Ask Chris B! \\
tcp\_accept\_count & DOUBLE & Number of TCP Accept for this process \\
tcp\_connect\_count & DOUBLE & Number of TCP Connects for this process \\
tcp\_disconnect\_count & DOUBLE & Number of TCP Disconnects for this process \\
tcp\_reconnect\_count & DOUBLE & Number of TCP Reconnects for this process \\
tcp\_recv\_count & DOUBLE & Number of TCP packets recevied by the process \\
tcp\_recv\_size & DOUBLE & Number of TCP bytes received by the process \\
tcp\_retransmit\_count & DOUBLE & Number of TCP retransmits by the process \\
tcp\_rs\_total & DOUBLE & Total TCP bytes sent/received by the process \\
tcp\_send\_count & DOUBLE & Number of TCP packets sent by the process \\
tcp\_send\_size & DOUBLE & Number of TCP bytes sent by the process \\
tcp\_send\_vs\_recv & DOUBLE & Ratio of TCP bytes sent vs received by the process \\
tcp\_tcpcopy\_count & DOUBLE & Number of TCP copy events during the process \\
tcp\_tcpcopy\_size & DOUBLE & TCP copy bytes \\
udp\_recv\_count & DOUBLE & Number of UDP packets received by the process \\
udp\_recv\_size & DOUBLE & Number of UDP bytes received by the process \\
udp\_rs\_total & DOUBLE & Total UDP bytes received by the process \\
udp\_send\_count & DOUBLE & Number of UDP packets sent by the process \\
udp\_send\_size & DOUBLE & Number of UDP bytes sent by the process \\
udp\_send\_vs\_recv & DOUBLE & Ratio of UDP bytes sent vs received by the process \\
\hline
\end{longtable}

\needspace{4\baselineskip}
\subsection{process\_path}
\label{sec:process_path}

For each process, has the path to its root (parent process).

\begin{longtable}{|p{5cm}|p{2cm}|p{9cm}|}
\hline
\textbf{Column} & \textbf{Data Type} & \textbf{Description} \\
\hline
\endfirsthead

\multicolumn{2}{c}%
{\tablename\ \thetable\ -- Continued from previous page} \\
\hline
\textbf{Column} & \textbf{Data Type} & \textbf{Description} \\
\hline
\endhead

\hline
\multicolumn{2}{r}{Continued on next page} \\
\endfoot

\hline
\endlastfoot

agent\_id & STRING & Agent ID is a unique id for a Wintap install on a host. It is created on initial startup after install. \\
hostname & STRING & Hostname data collected from \\
level & INTEGER & Distance from the kernel process \\
max\_level & INTEGER & Ignore: The same value level. Artifact of the processing to generate the paths. \\
os\_pid & INTEGER & Operating system process id \\
parent\_os\_pid & INTEGER & OS Process ID of parent process \\
parent\_pid\_hash & STRING & PID\_HASH of parent process \\
pid\_hash & STRING & Globally unique process hash (hash from hostname, os pid, and start\_time) \\
process\_name & STRING & Filename of executable \\
process\_path & STRING & Fully qualified path of the process executable \\
ptree & STRING & Path to kernel using only PROCESS\_NAMEs. Order and format is: =process->parent process->...  Ex: =winlogon.exe->smss.exe->smss.exe->ntoskrnl.exe \\
ptree\_list & STRING[] & Path to kernel using PID\_HASHes. Stored as an list. List is order from process to kernel process \\
ptree\_list\_tuples & STRUCT & Path to kernel using named tuples of PID\_HASH and PROCESS\_NAME (list of maps). List is order from process to kernel process \\
seq & INTEGER & Ignore: Artifact left over from process. Level+1. \\
\hline
\end{longtable}

\needspace{4\baselineskip}
\subsection{process\_registry}
\label{sec:process_registry}

Registry activity events. These are aggregated to PROCESS + Registry Key/Value

\begin{longtable}{|p{5cm}|p{2cm}|p{9cm}|}
\hline
\textbf{Column} & \textbf{Data Type} & \textbf{Description} \\
\hline
\endfirsthead

\multicolumn{2}{c}%
{\tablename\ \thetable\ -- Continued from previous page} \\
\hline
\textbf{Column} & \textbf{Data Type} & \textbf{Description} \\
\hline
\endhead

\hline
\multicolumn{2}{r}{Continued on next page} \\
\endfoot

\hline
\endlastfoot

activity\_type & STRING & Type of registry activity. Ex: CREATEKEY, READ, WRITE, DELETEVALUE, DELETEKEY \\
agent\_id & STRING & Agent ID is a unique id for a Wintap install on a host. It is created on initial startup after install. \\
event\_count & DOUBLE & Total ETW events \\
first\_seen & TIMESTAMP & Earliest time seen \\
hostname & STRING & Hostname data collected from \\
last\_seen & TIMESTAMP & Latest time seen \\
max\_event & TIMESTAMP & Maximum value in the interval \\
min\_event & TIMESTAMP & Minimum value in the interval \\
num\_raw\_rows & INT64 & Number of unique values in the related column. Typically, the related column should only have 1 value. \\
pid\_hash & STRING & Globally unique process hash (hash from hostname, os pid, and start\_time) \\
process\_name & STRING & Filename of executable \\
reg\_data & STRING & Registry data, which is the most detailed part of the registry key. Strangely, it isn't data \\
reg\_path & STRING & Registry path \\
reg\_value & STRING & Registry value read or written \\
\hline
\end{longtable}

\needspace{4\baselineskip}
\subsection{process\_registry\_summary}
\label{sec:process_registry_summary}

All registry activity events summarized to the PROCESS.

\begin{longtable}{|p{5cm}|p{2cm}|p{9cm}|}
\hline
\textbf{Column} & \textbf{Data Type} & \textbf{Description} \\
\hline
\endfirsthead

\multicolumn{2}{c}%
{\tablename\ \thetable\ -- Continued from previous page} \\
\hline
\textbf{Column} & \textbf{Data Type} & \textbf{Description} \\
\hline
\endhead

\hline
\multicolumn{2}{r}{Continued on next page} \\
\endfoot

\hline
\endlastfoot

agent\_id & STRING & Agent ID is a unique id for a Wintap install on a host. It is created on initial startup after install. \\
createkeys & DOUBLE & Number of createkeys by process \\
deletekeys & DOUBLE & Number of deletekeys by process \\
deletevalues & DOUBLE & Number of deletevalues by process \\
first\_seen & TIMESTAMP & Earliest time seen \\
hostname & STRING & Hostname data collected from \\
last\_seen & TIMESTAMP & Latest time seen \\
pid\_hash & STRING & Globally unique process hash (hash from hostname, os pid, and start\_time) \\
process\_name & STRING & Filename of executable \\
reads & DOUBLE & Number of reads by process \\
total\_activity\_types & DOUBLE & Total number of ETW events for all registry activity for a process \\
writes & DOUBLE & Number of writes by process \\
\hline
\end{longtable}

\needspace{4\baselineskip}
\subsection{process\_summary}
\label{sec:process_summary}

Provides a unified summary of all process-related activities that are directly collected by Wintap. This includes:n  Process metadata (e.g., process\_name, hostname, pid\_hash).n  Registry activity (e.g., reads, writes, key creation/deletion).n  File activity (e.g., file reads, writes, and other events).n  Network activity (e.g., connections, data sent/received).n  Image loads (e.g., DLLs loaded by the process).n  Host-level metadata (e.g., OS, architecture).

\begin{longtable}{|p{5cm}|p{2cm}|p{9cm}|}
\hline
\textbf{Column} & \textbf{Data Type} & \textbf{Description} \\
\hline
\endfirsthead

\multicolumn{2}{c}%
{\tablename\ \thetable\ -- Continued from previous page} \\
\hline
\textbf{Column} & \textbf{Data Type} & \textbf{Description} \\
\hline
\endhead

\hline
\multicolumn{2}{r}{Continued on next page} \\
\endfoot

\hline
\endlastfoot

Close\_Events & DOUBLE & Number of file close events \\
Create\_Events & DOUBLE & Number of file create events \\
Delete\_Events & DOUBLE & Number of file delete events \\
Read\_Bytes & DOUBLE & File bytes read \\
Read\_Events & DOUBLE & Number of file read events \\
Rename\_Events & DOUBLE & Number of file rename events \\
SetInfo\_Events & DOUBLE &  \\
Write\_Bytes & DOUBLE & File bytes written \\
Write\_Events & DOUBLE & Number of file write events \\
agent\_id & STRING & Agent ID is a unique id for a Wintap install on a host. It is created on initial startup after install. \\
arch & STRING & Host architecture \\
args & STRING & Process arguments. Note: as a simple string, no parsing done \\
avg\_bytes & DOUBLE & Average bytes over the interval \\
avg\_packets & DOUBLE & Average packets over the interval \\
commit\_charge & INT64 & Commit charge is the total amount of virtual memory guaranteed for all processes to fit in physical memory and the page file. \\
commit\_peak & INT64 & Total memory commited \\
conn\_id\_count & INT64 & Number of unique network connection IDs (5-tuples) \\
cpu\_cycle\_count & INT64 & Always zero (note: more than zero cycles were generally used in a process) \\
cpu\_utilization & INTEGER & Ignore: unable to collect. This *should* be a value indicating the percentage of cpu utilized by the process over its life, but turns out be harder to collect than expected. \\
dll\_first\_seen & TIMESTAMP & Earliest time seen \\
dll\_last\_seen & TIMESTAMP & Latest time seen \\
dll\_num\_uniq\_files & INT64 & Number of unique DDLs by name \\
dlls & STRING[] & List of DDL names as an array \\
duration\_seconds & DOUBLE & Elapsed process execution in seconds \\
exit\_code & INT64 & Process exit code \\
file\_first\_seen & TIMESTAMP & Earliest time seen \\
file\_id & STRING & Unique ID for a file. Hash of hostname+fully qualified filename \\
file\_last\_seen & TIMESTAMP & Latest time seen \\
file\_md5 & STRING & MD5 hash of the binary process/dll \\
file\_num\_raw\_rows & DOUBLE & Number of rows of file activity from the raw data. \\
file\_sha2 & STRING & SHA2 hash of the file contents \\
filename & STRING & Fully qualified filename \\
first\_seen & TIMESTAMP & Earliest time seen \\
hard\_fault\_count & INTEGER & Count of hard page faults \\
hostname & STRING & Hostname data collected from \\
last\_seen & TIMESTAMP & Latest time seen \\
max\_bytes & DOUBLE & Max bytes in the interval \\
max\_packets & DOUBLE & Max packets in the interval \\
min\_bytes & DOUBLE & Min bytes in the interval \\
min\_packets & DOUBLE & Max packets in the interval \\
net\_first\_seen & TIMESTAMP & Earliest time seen \\
net\_last\_seen & TIMESTAMP & Latest time seen \\
net\_num\_raw\_rows & DOUBLE & Number of rows in the raw data \\
net\_recv\_size & DOUBLE & Total bytes received in the interval \\
net\_rs\_total & DOUBLE & Total bytes sent/received in the interval \\
net\_send\_size & DOUBLE & Total bytes sent in the interval \\
net\_send\_vs\_recv & DOUBLE & Ratio of bytes sent/received in the interval \\
net\_total\_events & DOUBLE & Total network events in the interval \\
net\_total\_size & DOUBLE & Total bytes (send/receive) in the interval \\
num\_agent\_id & INT64 & Number of unique values in the related column. Typically, the related column should only have 1 value. \\
num\_args & INT64 & Number of unique values in the related column. Typically, the related column should only have 1 value. \\
num\_file\_md5 & INT64 & Number of unique values in the related column. Typically, the related column should only have 1 value. \\
num\_file\_sha2 & INT64 & Number of unique values in the related column. Typically, the related column should only have 1 value. \\
num\_null\_filename & DOUBLE & Number of unique values in the related column. Typically, the related column should only have 1 value. \\
num\_parent\_os\_pid & INT64 & Number of unique values in the related column. Typically, the related column should only have 1 value. \\
num\_parent\_pid\_hash & INT64 & Number of unique values in the related column. Typically, the related column should only have 1 value. \\
num\_process\_name & INT64 & Number of unique values in the related column. Typically, the related column should only have 1 value. \\
num\_process\_path & INT64 & Number of unique values in the related column. Typically, the related column should only have 1 value. \\
num\_process\_start & DOUBLE & Number of unique values in the related column. Typically, the related column should only have 1 value. \\
num\_process\_stop & DOUBLE & Number of unique values in the related column. Typically, the related column should only have 1 value. \\
num\_uniq\_file\_hash & INT64 & Number of unique values in the related column. Typically, the related column should only have 1 value. \\
num\_user\_name & INT64 & Number of unique values in the related column. Typically, the related column should only have 1 value. \\
os & STRING & Operating System \\
os\_family & STRING & OS Family: windows, linux, osx \\
os\_pid & INTEGER & Operating system process id \\
os\_version & STRING & Operating System Version \\
parent\_os\_pid & INTEGER & OS Process ID of parent process \\
parent\_pid\_hash & STRING & PID\_HASH of parent process \\
pid\_hash & STRING & Globally unique process hash (hash from hostname, os pid, and start\_time) \\
process\_name & STRING & Filename of executable \\
process\_path & STRING & Fully qualified path of the process executable \\
process\_started & TIMESTAMP & Process start timestamp \\
process\_started\_seconds & DOUBLE & Process start time in unix epoch seconds \\
process\_stop\_seconds & DOUBLE & Process stop in unix epoch seconds \\
process\_term & TIMESTAMP & Process termination timestamp \\
read\_operation\_count & INT64 & Count of number of read operations \\
read\_transfer\_kilobytes & INT64 & Number of K bytes read \\
reg\_createkeys & DOUBLE & Registry keys created \\
reg\_deletekeys & DOUBLE & Registry keys deleted \\
reg\_deletevalues & DOUBLE & Registry values deleted \\
reg\_first\_seen & TIMESTAMP & Earliest time seen \\
reg\_last\_seen & TIMESTAMP & Latest time seen \\
reg\_reads & DOUBLE & Registry read events \\
reg\_totals & DOUBLE & Total registry events \\
reg\_writes & DOUBLE & Registry write events \\
sq\_size & DOUBLE & Squared size of bytes in network packet. Ask Chris B! \\
tcp\_accept\_count & DOUBLE & TCP Accepts \\
tcp\_connect\_count & DOUBLE & TCP connects \\
tcp\_disconnect\_count & DOUBLE & TCP disconnects \\
tcp\_reconnect\_count & DOUBLE & TCP reconnects \\
tcp\_recv\_count & DOUBLE & TCP packets received \\
tcp\_recv\_size & DOUBLE & TCP bytes received \\
tcp\_retransmit\_count & DOUBLE & TCP packets retransmitted \\
tcp\_rs\_total & DOUBLE & TCP bytes total (send/received) \\
tcp\_send\_count & DOUBLE & TCP bytes sent \\
tcp\_send\_size & DOUBLE & TCP packets sent \\
tcp\_send\_vs\_recv & DOUBLE & Ratio of TCP bytes sent/received \\
tcp\_tcpcopy\_count & DOUBLE & TCP copy events. We think this is an event that is internal to the network stack on the host. \\
tcp\_tcpcopy\_size & DOUBLE & TCP copy bytes \\
token\_elevation\_type & INTEGER & Indicates UAC privilege status: 1 (Default) = no linked token, UAC disabled or non-admin account; 2 (Full) = elevated administrator with full rights after UAC approval; 3 (Limited) = filtered admin token before elevation with available full token; 4 = undocumented value, treat as unknown/reserved. \\
udp\_recv\_count & DOUBLE & UDP bytes received \\
udp\_recv\_size & DOUBLE & UDP events received \\
udp\_rs\_total & DOUBLE & Total UDP bytes sent/received \\
udp\_send\_count & DOUBLE & UDP packets sent \\
udp\_send\_size & DOUBLE & UDP bytes sent \\
udp\_send\_vs\_recv & DOUBLE & Ratio of UDP bytes sent/recieved \\
user\_name & STRING & Username. Note that on windows there are many usernames that represent system or other background activity \\
write\_operation\_count & INT64 & Count of number of write operations \\
write\_transfer\_kilobytes & INT64 & Number of K bytes written \\
\hline
\end{longtable}

\needspace{4\baselineskip}
\subsection{process\_uber\_summary}
\label{sec:process_uber_summary}

The process\_uber\_summary view is an enhanced version of the process\_summary view, incorporating additional threat intelligence and labeling data. It combines process activity summaries with external threat indicators such as SIGMA rules, MITRE ATT\&CK techniques, LOLBAS (Living Off the Land Binaries and Scripts), and NetworkX graph labels. This comprehensive view provides a detailed and enriched dataset for advanced threat analysis and detection.

\begin{longtable}{|p{5cm}|p{2cm}|p{9cm}|}
\hline
\textbf{Column} & \textbf{Data Type} & \textbf{Description} \\
\hline
\endfirsthead

\multicolumn{2}{c}%
{\tablename\ \thetable\ -- Continued from previous page} \\
\hline
\textbf{Column} & \textbf{Data Type} & \textbf{Description} \\
\hline
\endhead

\hline
\multicolumn{2}{r}{Continued on next page} \\
\endfoot

\hline
\endlastfoot

Close\_Events & DOUBLE & Number of file close events \\
Create\_Events & DOUBLE & Number of file create events \\
Delete\_Events & DOUBLE & Number of file delete events \\
Read\_Bytes & DOUBLE & File bytes read \\
Read\_Events & DOUBLE & Number of file read events \\
Rename\_Events & DOUBLE & Number of file rename events \\
SetInfo\_Events & DOUBLE &  \\
Write\_Bytes & DOUBLE & File bytes written \\
Write\_Events & DOUBLE & Number of file write events \\
agent\_id & STRING & Agent ID is a unique id for a Wintap install on a host. It is created on initial startup after install. \\
arch & STRING & Host architecture \\
args & STRING & Process arguments. Note: as a simple string, no parsing done \\
avg\_bytes & DOUBLE & Average bytes over the interval \\
avg\_packets & DOUBLE & Average packets over the interval \\
commit\_charge & INT64 & Commit charge is the total amount of virtual memory guaranteed for all processes to fit in physical memory and the page file. \\
commit\_peak & INT64 & Total memory commited \\
conn\_id\_count & INT64 & Number of unique network connection IDs (5-tuples) \\
cpu\_cycle\_count & INT64 & Always zero (note: more than zero cycles were generally used in a process) \\
cpu\_utilization & INTEGER & Ignore: unable to collect. This *should* be a value indicating the percentage of cpu utilized by the process over its life, but turns out be harder to collect than expected. \\
critical\_num\_sigma\_hits & DOUBLE & Number of unique Sigma rules hit \\
dll\_first\_seen & TIMESTAMP & Earliest time seen \\
dll\_last\_seen & TIMESTAMP & Latest time seen \\
dll\_num\_uniq\_files & INT64 & Number of unique DDLs by name \\
dlls & STRING[] & List of DDL names as an array \\
duration\_seconds & DOUBLE & Elapsed process execution in seconds \\
exit\_code & INT64 & Process exit code \\
file\_first\_seen & TIMESTAMP & Earliest time seen \\
file\_id & STRING & Unique ID for a file. Hash of hostname+fully qualified filename \\
file\_last\_seen & TIMESTAMP & Latest time seen \\
file\_md5 & STRING & MD5 hash of the binary process/dll \\
file\_num\_raw\_rows & DOUBLE & Number of rows of file activity from the raw data. \\
file\_sha2 & STRING & SHA2 hash of the file contents \\
filename & STRING & Fully qualified filename \\
first\_seen & TIMESTAMP & Earliest time seen \\
hard\_fault\_count & INTEGER & Count of hard page faults \\
high\_num\_sigma\_hits & INT64 & Number of unique Sigma rules hit \\
high\_num\_sigma\_rows & INT64 & Number of times Sigma rules hit for this entity \\
hostname & STRING & Hostname data collected from \\
label\_num\_hits & INT64 & Number of unique labels hit \\
label\_num\_sources & INT64 & Number of unique sources of labels \\
label\_num\_uniq\_annotations & INT64 & Number of unqiue annotations \\
label\_source & STRING & For now, only a single source: "networkx" \\
last\_seen & TIMESTAMP & Latest time seen \\
lolbas\_cats & STRING[] & List categories from hits \\
lolbas\_mitre & STRING[] & List of Mitre codes from hits \\
lolbas\_num\_rows & INT64 & Number of hits \\
lolbas\_privs & STRING[] & List of priviledges from hist \\
low\_num\_sigma\_hits & INT64 & Number of unique Sigma rules hit \\
low\_num\_sigma\_rows & INT64 & Number of times Sigma rules hit for this entity \\
max\_bytes & DOUBLE & Max bytes in the interval \\
max\_packets & DOUBLE & Max packets in the interval \\
medium\_num\_sigma\_hits & INT64 & Number of unique Sigma rules hit \\
medium\_num\_sigma\_rows & INT64 & Number of times Sigma rules hit for this entity \\
min\_bytes & DOUBLE & Min bytes in the interval \\
min\_packets & DOUBLE & Max packets in the interval \\
mitre\_analytic\_ids & STRING[] & List of analytic IDs from hits \\
mitre\_analytic\_types & STRING[][] & List of ananlytic types from hits \\
mitre\_information\_domains & STRING[] & List domains from hits. Ex: Analytic, Host, Network \\
mitre\_num\_rows & INT64 & Number of hits \\
mitre\_subtypes & STRING[][] & List of subtypes from hits. Ex: Map building, Anomaly, Hostflow.  Process. \\
net\_first\_seen & TIMESTAMP & Earliest time seen \\
net\_last\_seen & TIMESTAMP & Latest time seen \\
net\_num\_raw\_rows & DOUBLE & Number of rows in the raw data \\
net\_recv\_size & DOUBLE & Total bytes received in the interval \\
net\_rs\_total & DOUBLE & Total bytes sent/received in the interval \\
net\_send\_size & DOUBLE & Total bytes sent in the interval \\
net\_send\_vs\_recv & DOUBLE & Ratio of bytes sent/received in the interval \\
net\_total\_events & DOUBLE & Total network events in the interval \\
net\_total\_size & DOUBLE & Total bytes (send/receive) in the interval \\
num\_agent\_id & INT64 & Number of unique values in the related column. Typically, the related column should only have 1 value. \\
num\_args & INT64 & Number of unique values in the related column. Typically, the related column should only have 1 value. \\
num\_file\_md5 & INT64 & Number of unique values in the related column. Typically, the related column should only have 1 value. \\
num\_file\_sha2 & INT64 & Number of unique values in the related column. Typically, the related column should only have 1 value. \\
num\_null\_filename & DOUBLE & Number of unique values in the related column. Typically, the related column should only have 1 value. \\
num\_parent\_os\_pid & INT64 & Number of unique values in the related column. Typically, the related column should only have 1 value. \\
num\_parent\_pid\_hash & INT64 & Number of unique values in the related column. Typically, the related column should only have 1 value. \\
num\_process\_name & INT64 & Number of unique values in the related column. Typically, the related column should only have 1 value. \\
num\_process\_path & INT64 & Number of unique values in the related column. Typically, the related column should only have 1 value. \\
num\_process\_start & DOUBLE & Number of unique values in the related column. Typically, the related column should only have 1 value. \\
num\_process\_stop & DOUBLE & Number of unique values in the related column. Typically, the related column should only have 1 value. \\
num\_uniq\_file\_hash & INT64 & Number of unique values in the related column. Typically, the related column should only have 1 value. \\
num\_user\_name & INT64 & Number of unique values in the related column. Typically, the related column should only have 1 value. \\
os & STRING & Operating System \\
os\_family & STRING & OS Family: windows, linux, osx \\
os\_pid & INTEGER & Operating system process id \\
os\_version & STRING & Operating System Version \\
parent\_os\_pid & INTEGER & OS Process ID of parent process \\
parent\_pid\_hash & STRING & PID\_HASH of parent process \\
pid\_hash & STRING & Globally unique process hash (hash from hostname, os pid, and start\_time) \\
process\_name & STRING & Filename of executable \\
process\_path & STRING & Fully qualified path of the process executable \\
process\_started & TIMESTAMP & Process start timestamp \\
process\_started\_seconds & DOUBLE & Process start time in unix epoch seconds \\
process\_stop\_seconds & DOUBLE & Process stop in unix epoch seconds \\
process\_term & TIMESTAMP & Process termination timestamp \\
read\_operation\_count & INT64 & Count of number of read operations \\
read\_transfer\_kilobytes & INT64 & Number of K bytes read \\
reg\_createkeys & DOUBLE & Registry keys created \\
reg\_deletekeys & DOUBLE & Registry keys deleted \\
reg\_deletevalues & DOUBLE & Registry values deleted \\
reg\_first\_seen & TIMESTAMP & Earliest time seen \\
reg\_last\_seen & TIMESTAMP & Latest time seen \\
reg\_reads & DOUBLE & Registry read events \\
reg\_totals & DOUBLE & Total registry events \\
reg\_writes & DOUBLE & Registry write events \\
sq\_size & DOUBLE & Squared size of bytes in network packet. Ask Chris B! \\
tcp\_accept\_count & DOUBLE & TCP Accepts \\
tcp\_connect\_count & DOUBLE & TCP connects \\
tcp\_disconnect\_count & DOUBLE & TCP disconnects \\
tcp\_reconnect\_count & DOUBLE & TCP reconnects \\
tcp\_recv\_count & DOUBLE & TCP packets received \\
tcp\_recv\_size & DOUBLE & TCP bytes received \\
tcp\_retransmit\_count & DOUBLE & TCP packets retransmitted \\
tcp\_rs\_total & DOUBLE & TCP bytes total (send/received) \\
tcp\_send\_count & DOUBLE & TCP bytes sent \\
tcp\_send\_size & DOUBLE & TCP packets sent \\
tcp\_send\_vs\_recv & DOUBLE & Ratio of TCP bytes sent/received \\
tcp\_tcpcopy\_count & DOUBLE & TCP copy events. We think this is an event that is internal to the network stack on the host. \\
tcp\_tcpcopy\_size & DOUBLE & TCP copy bytes \\
token\_elevation\_type & INTEGER & Indicates UAC privilege status: 1 (Default) = no linked token, UAC disabled or non-admin account; 2 (Full) = elevated administrator with full rights after UAC approval; 3 (Limited) = filtered admin token before elevation with available full token; 4 = undocumented value, treat as unknown/reserved. \\
total\_sigma\_hits & DOUBLE & Total Sigma hits over all crticality levels \\
udp\_recv\_count & DOUBLE & UDP bytes received \\
udp\_recv\_size & DOUBLE & UDP events received \\
udp\_rs\_total & DOUBLE & Total UDP bytes sent/received \\
udp\_send\_count & DOUBLE & UDP packets sent \\
udp\_send\_size & DOUBLE & UDP bytes sent \\
udp\_send\_vs\_recv & DOUBLE & Ratio of UDP bytes sent/recieved \\
user\_name & STRING & Username. Note that on windows there are many usernames that represent system or other background activity \\
write\_operation\_count & INT64 & Count of number of write operations \\
write\_transfer\_kilobytes & INT64 & Number of K bytes written \\
\hline
\end{longtable}

\needspace{4\baselineskip}
\subsection{sigma\_labels}
\label{sec:sigma_labels}

The sigma\_labels table is a key component in the threat detection pipeline, as it links process or entity activity to specific SIGMA rules. SIGMA is a standardized rule format for describing log-based detection patterns, often used for identifying suspicious or malicious behavior in systems. The sigma\_labels table provides a mapping between entities (like processes) and the SIGMA rules that matched their activity, along with metadata about the severity and type of the detection.

\begin{longtable}{|p{5cm}|p{2cm}|p{9cm}|}
\hline
\textbf{Column} & \textbf{Data Type} & \textbf{Description} \\
\hline
\endfirsthead

\multicolumn{2}{c}%
{\tablename\ \thetable\ -- Continued from previous page} \\
\hline
\textbf{Column} & \textbf{Data Type} & \textbf{Description} \\
\hline
\endhead

\hline
\multicolumn{2}{r}{Continued on next page} \\
\endfoot

\hline
\endlastfoot

analytic\_id & STRING & Specific SIGMA rule hit. Value is a GUID \\
entity & STRING & Wintap entity ID (hash). The type of entity (PID\_HASH, CONN\_ID, etc) is defined in the ENTITY\_TYPE column. Note: Currently, only PID\_HASHes are identified. \\
entity\_type & STRING & Wintap entity type. Determines what kind of hash is in the ENTITY column. Ex: PID\_HASH, CONN\_ID, FILE\_ID \\
time & TIMESTAMP & Timestamp of activity \\
\hline
\end{longtable}

\needspace{4\baselineskip}
\subsection{sigma\_labels\_summary}
\label{sec:sigma_labels_summary}

Aggregates data from sigma\_labels by PID\_HASH and pivots the severity levels into separate columns (e.g., critical\_num\_sigma\_hits, high\_num\_sigma\_hits).

\begin{longtable}{|p{5cm}|p{2cm}|p{9cm}|}
\hline
\textbf{Column} & \textbf{Data Type} & \textbf{Description} \\
\hline
\endfirsthead

\multicolumn{2}{c}%
{\tablename\ \thetable\ -- Continued from previous page} \\
\hline
\textbf{Column} & \textbf{Data Type} & \textbf{Description} \\
\hline
\endhead

\hline
\multicolumn{2}{r}{Continued on next page} \\
\endfoot

\hline
\endlastfoot

critical\_num\_sigma\_hits & DOUBLE & Number of unique Sigma rules hit \\
high\_num\_sigma\_hits & INT64 & Number of unique Sigma rules hit \\
high\_num\_sigma\_rows & INT64 & Number of times Sigma rules hit for this entity \\
low\_num\_sigma\_hits & INT64 & Number of unique Sigma rules hit \\
low\_num\_sigma\_rows & INT64 & Number of times Sigma rules hit for this entity \\
medium\_num\_sigma\_hits & INT64 & Number of unique Sigma rules hit \\
medium\_num\_sigma\_rows & INT64 & Number of times Sigma rules hit for this entity \\
pid\_hash & STRING & Globally unique process hash (hash from hostname, os pid, and start\_time) \\
\hline
\end{longtable}

\twocolumn\end{document}